\documentclass[11pt]{article}

% ============================================================================
% HISTORICAL / EXPLORATORY DOCUMENT
% ============================================================================
% This detailed Part-7 file records the pre-freeze actuator-realization
% investigation, including affine/hybrid candidates and open C2.8/C2.9 ideas.
% It is NOT the authoritative final actuator contract.
%
% Authoritative final contract:
%   docs/mathematics/contracts/
%   PINODE_EPSR_Part7_MPC_Actuator_Realization_Simplified_FROZEN.tex
%
% Repository ownership:
%   Clean LaTeX/docs live under the EPSR code repo.
%   Scientific EnergyPlus experiment evidence lives under the separate data root.
% ============================================================================

\usepackage[margin=1in]{geometry}
\usepackage{amsmath,amssymb,mathtools}
\usepackage{booktabs,longtable,array}
\usepackage{enumitem}
\usepackage{microtype}
\usepackage[hidelinks]{hyperref}
\usepackage{xcolor}

\newcommand{\mdot}{\dot m}
\newcommand{\Qx}{Q_{\mathrm{HVAC},X}}
\newcommand{\Qac}{Q_{\mathrm{AC}}}
\newcommand{\Phvac}{P_{\mathrm{HVAC}}}
\newcommand{\uStar}{u^\star}
\newcommand{\uEP}{u^{\mathrm{EP}}}
\newcommand{\aEP}{a^{\mathrm{EP}}}
\newcommand{\R}{\mathbb{R}}
\newcommand{\Frozen}{\textbf{[EMPIRICALLY FROZEN]}}
\newcommand{\Candidate}{\textbf{[CURRENT CANDIDATE]}}
\newcommand{\OpenItem}{\textbf{[OPEN / TO IDENTIFY]}}
\newcommand{\PaperRoot}{\texttt{Paper\_PINODE\_EPSR}}

\title{PINODE/EPSR MPC and EnergyPlus Actuator Realization\\
Detailed Mathematical Contract and Experimental Identification Record}
\author{ScaleBridge / PINODE--EPSR Working Mathematical Contract}
\date{29 August 2026 --- Version 5 (compile-clean)}

\begin{document}
\maketitle

\begin{abstract}
This document extends the existing PINODE/EPSR MPC mathematics with the
EnergyPlus actuator-realization layer required for closed-loop experiments on
the RestaurantFastFood/Buffalo prototype.  The upper-level MPC retains the
physically meaningful control variables supply-air mass flow
$\mdot$ and supply-air temperature $T_{\mathrm{sa}}$.  A separate lower-level realization
map converts those physical requests into the low-level EnergyPlus commands
actually accepted by the two packaged single-zone (PSZ) HVAC systems serving
Dining and Kitchen.  The formulation deliberately separates (i) desired MPC
inputs, (ii) EnergyPlus actuator commands, and (iii) physically realized HVAC
inputs.  Numerical experiments C1--C2.4 are incorporated explicitly, including
the exact identified fan-flow inverse and the current evidence for
mode-dependent thermal realization.  Items that are already established by
experiment are marked separately from models that remain to be identified.
\end{abstract}

\section{Purpose and relationship to the existing Part-6 MPC contract}

The existing PINODE/EPSR MPC contract keeps all four thermal-model families in
continuous time and advances them over a fixed control interval using numerical
integration.  The closed-loop control interval is
\begin{equation}
\Delta t_{\mathrm{MPC}} = 300~\mathrm{s}.
\label{eq:mpc-dt}
\end{equation}
The intended physical HVAC path in the existing contract is
\begin{equation}
(\mdot,T_{\mathrm{sa}},T_z)
\longrightarrow
\Qx
\longrightarrow
G_Q
\longrightarrow
\Qac
\longrightarrow
\dot x = F_\theta(x,\Qac,d)
\longrightarrow
x_{k+1},
\label{eq:existing-thermal-chain}
\end{equation}
while the electrical-power path remains separate,
\begin{equation}
|\Qac|
\longrightarrow
G_P
\longrightarrow
\Phvac.
\label{eq:existing-power-chain}
\end{equation}
Here $G_Q$ is the common Phase-C signed thermal-input mapping and $G_P$ is the
common Phase-C HVAC electrical-power mapping.  In particular, $\Phvac$ is an
MPC objective/output quantity and does not enter the thermal ODE right-hand
side.

The present document adds a realization layer \emph{before}
Eq.~\eqref{eq:existing-thermal-chain}.  It answers the practical question:
\begin{quote}
Given the MPC-requested physical pair $(\mdot^\star,T_{\mathrm{sa}}^\star)$, what
low-level EnergyPlus actuator commands must be issued so that the actual
RestaurantFastFood HVAC system realizes that pair as closely as physically
possible?
\end{quote}

\section{System under study}

The current controlled EnergyPlus plant is the RestaurantFastFood prototype in
Buffalo, using the SHA-verified EnergyPlus 24.1 transition candidate
\begin{verbatim}
RestaurantFastFood_Buffalo_EPlus24.1_FINAL_CANDIDATE.idf

SHA256:
da86d8c3c782424d2d4f30e29a97220c9b9042f721e83591f43762f34df77101
\end{verbatim}

The physical thermal zones of interest are
\begin{equation}
\mathcal Z = \{D,K\},
\end{equation}
where $D$ denotes Dining and $K$ denotes Kitchen.

Both zones are served by independent
\texttt{AirLoopHVAC:UnitarySystem} packaged single-zone systems with
\texttt{Control Type = Load}.  The major physical chains are:

\begin{verbatim}
DINING: PSZ-AC_1:1

return / supply-equipment inlet
          |
          +---- outdoor air
          |
          v
      OA mixer
          |
          v
  DX cooling coil
          |
          v
    heating coil
          |
          v
   Fan:SystemModel
          |
          v
 PSZ supply outlet
          |
          v
 Dining Direct Air
          |
          v
      DINING ZONE


KITCHEN: PSZ-AC_2:2

return / supply-equipment inlet
          |
          +---- outdoor / makeup air
          |
          v
      OA mixer
          |
          v
  DX cooling coil
          |
          v
    heating coil
          |
          v
   Fan:SystemModel
          |
          v
 PSZ supply outlet
          |
          v
 Kitchen Direct Air
          |
          v
      KITCHEN ZONE
\end{verbatim}

The Kitchen system has a particularly important outdoor/makeup-air burden.
Consequently, positive heating-coil output does not necessarily imply positive
sensible heat delivery relative to the Kitchen zone temperature.

\section{Notation registry}

\begin{longtable}{>{\raggedright\arraybackslash}p{0.20\textwidth}
                  >{\raggedright\arraybackslash}p{0.58\textwidth}
                  >{\raggedright\arraybackslash}p{0.14\textwidth}}
\toprule
Symbol & Meaning & Units \\
\midrule
\endfirsthead
\toprule
Symbol & Meaning & Units \\
\midrule
\endhead
$i\in\mathcal Z$ & Zone index; $D$ = Dining, $K$ = Kitchen & -- \\
$k$ & MPC control-step index & -- \\
$j$ & Prediction-horizon index & -- \\
$\Delta t_{\mathrm{MPC}}$ & Closed-loop control interval & s \\
$x_{i,k}$ & Thermal-model state for zone/model block $i$ & model-dependent \\
$\hat x_{i,k}$ & Observer/encoder state estimate used by MPC & model-dependent \\
$d_{i,k}$ & Exogenous disturbances/forecasts & model-dependent \\
$\xi_{i,k}$ & Local HVAC realization state/context & mixed \\
$T_{z,i,k}$ & Zone air temperature & $^\circ$C \\
$\mdot_{i,k}^\star$ & Physical supply-air mass-flow request produced by MPC & kg/s \\
$T_{\mathrm{sa},i,k}^\star$ & Physical supply-air-temperature request produced by MPC & $^\circ$C \\
$u_{i,k}^\star$ & Desired physical MPC HVAC input $[\mdot_i^\star,T_{\mathrm{sa},i}^\star]^\top$ & mixed \\
$a_{m,i,k}$ & Low-level EnergyPlus fan mass-flow actuator command & kg/s \\
$a_{q,i,k}$ & Low-level EnergyPlus unitary sensible-load request & W \\
$a_{s,i,k}$ & Low-level EnergyPlus DX-speed command when used & dimensionless \\
$a_{i,k}$ & Complete low-level EnergyPlus actuator vector & mixed \\
$\mdot_{i,k}^{\mathrm{EP}}$ & Realized supply/direct-air-node mass flow & kg/s \\
$T_{\mathrm{sa},i,k}^{\mathrm{EP}}$ & Realized supply/direct-air-node temperature & $^\circ$C \\
$u_{i,k}^{\mathrm{EP}}$ & Realized physical HVAC input & mixed \\
$c_p$ & Specific heat of air used consistently in control equations & J/(kg K) \\
$\Qx$ & Raw air-side sensible heat quantity prior to Phase-C $G_Q$ & W \\
$\Qac$ & Signed Phase-C thermal input supplied to the thermal ODE & W \\
$\Phvac$ & Phase-C HVAC electrical-power output & W \\
$G_Q$ & Common Phase-C signed thermal-input mapping & -- \\
$G_P$ & Common Phase-C HVAC electrical-power mapping & -- \\
$g_i$ & Forward low-level actuator-to-physical-HVAC realization map & -- \\
$f_i$ & Inverse/realizer map from physical HVAC target to low-level commands & -- \\
$\alpha_i,\beta_i$ & Local gain and offset for an affine thermal-realization model & --, W \\
$\mathcal A_i$ & Admissible low-level actuator-command set & mixed \\
$\mathcal U_i^{\mathrm{EP}}(\xi)$ & Physically realizable $(\mdot,T_{\mathrm{sa}})$ set under EnergyPlus state $\xi$ & mixed \\
$\mathcal U_i^{\mathrm{train}}$ & Region supported by thermal-model training data & mixed \\
$\mathcal U_i^{\mathrm{MPC}}$ & Final MPC physical-input feasible set & mixed \\
$e^u_i$ & Physical actuator-realization error & mixed \\
$e_i^{\mdot}$ & Mass-flow realization error & kg/s \\
$e_i^T$ & Supply-temperature realization error & K or $^\circ$C difference \\
\bottomrule
\end{longtable}

\paragraph{Specific-heat convention.}
The equations in this document retain the symbolic constant $c_p$ so one
numerical convention is used throughout the final paper and implementation.
The earlier Part-6 working contract used approximately
$c_p=1005~\mathrm{J/(kg\,K)}$, whereas the C2.3/C2.4 screening scripts used
$1006~\mathrm{J/(kg\,K)}$.  The final implementation should freeze one value
and use it consistently in the MPC, actuator realizer, Phase-C interface, and
all reported metrics.  This numerical difference is small but must not remain
implicit in the final paper.

\section{Two-layer control architecture}

The fundamental architecture is

\begin{verbatim}
 weather / gains / price / forecasts
                 |
                 v
        +------------------+
        |       MPC        |
        |                  |
        | outputs physical |
        |  m_dot* , T_sa*  |
        +---------+--------+
                  |
                  | u* = [m_dot*, T_sa*]
                  v
        +---------------------------+
        | actuator realization f_i  |
        |                           |
        | physical request          |
        |      ->                   |
        | EnergyPlus commands       |
        +-------------+-------------+
                      |
                      | [fan command,
                      |  load request,
                      |  DX command, ...]
                      v
        +---------------------------+
        |   EnergyPlus PSZ plant    |
        +-------------+-------------+
                      |
                      v
          m_dot_EP , T_sa_EP
                      |
                      +------> realized Q_HVAC
                      |
                      +------> zone / next MPC step
\end{verbatim}

For each zone,
\begin{equation}
u_{i,k}^\star
=
\begin{bmatrix}
\mdot_{i,k}^\star\\
T_{\mathrm{sa},i,k}^\star
\end{bmatrix}.
\label{eq:physical-mpc-input}
\end{equation}

The low-level command vector is
\begin{equation}
a_{i,k}
=
\begin{bmatrix}
a_{m,i,k}\\
a_{q,i,k}\\
a_{s,i,k}\\
\vdots
\end{bmatrix}.
\label{eq:low-level-actuator-vector}
\end{equation}

The true EnergyPlus plant defines a forward realization map
\begin{equation}
u_{i,k}^{\mathrm{EP}}
=
g_i(a_{i,k},\xi_{i,k})
=
\begin{bmatrix}
\mdot_{i,k}^{\mathrm{EP}}\\
T_{\mathrm{sa},i,k}^{\mathrm{EP}}
\end{bmatrix},
\label{eq:forward-realization-map}
\end{equation}
where $\xi_{i,k}$ may include zone state, outdoor/mixed-air state, humidity,
equipment mode, controller state, and other quantities required to explain the
local PSZ response.

The low-level controller seeks an inverse or realization policy
\begin{equation}
a_{i,k}
=
f_i
\left(
u_{i,k}^\star,\xi_{i,k}
\right)
=
f_i
\left(
\mdot_{i,k}^\star,
T_{\mathrm{sa},i,k}^\star,
\xi_{i,k}
\right).
\label{eq:inverse-realizer}
\end{equation}

The distinction between Eqs.~\eqref{eq:physical-mpc-input} and
\eqref{eq:low-level-actuator-vector} is essential: the MPC variables remain
physical variables even when the EnergyPlus realization requires different
low-level commands.

\section{Physical HVAC heat relation}

At the zone-supply interface, the desired raw sensible heat associated with the
MPC request is
\begin{equation}
Q_{X,i,k}^{\star}
=
\mdot_{i,k}^{\star}
c_p
\left(
T_{\mathrm{sa},i,k}^{\star}-T_{z,i,k}
\right).
\label{eq:desired-airside-heat}
\end{equation}

The realized quantity is
\begin{equation}
Q_{X,i,k}^{\mathrm{EP}}
=
\mdot_{i,k}^{\mathrm{EP}}
c_p
\left(
T_{\mathrm{sa},i,k}^{\mathrm{EP}}-T_{z,i,k}
\right).
\label{eq:realized-airside-heat}
\end{equation}

The existing paper pipeline then applies the common Phase-C mapping,
\begin{equation}
Q_{\mathrm{AC},i,k}
=
G_Q
\left(
Q_{X,i,k},\zeta_{Q,i,k}
\right),
\label{eq:phase-c-q-map}
\end{equation}
where $\zeta_Q$ denotes any additional regressors required by the frozen Phase-C
model.  The thermal ODE uses $\Qac$, not electrical power.

Electrical HVAC power is computed separately:
\begin{equation}
P_{\mathrm{HVAC},i,k}
=
G_P
\left(
|Q_{\mathrm{AC},i,k}|,\zeta_{P,i,k}
\right).
\label{eq:phase-c-p-map}
\end{equation}

Thus,
\begin{verbatim}
m_dot , T_sa , T_zone
          |
          v
 Q_HVAC,X = m_dot cp (T_sa - T_zone)
          |
          v
       Phase-C G_Q
          |
          v
        signed Q_AC
          |
          +-------------> continuous thermal model / RK4
          |
          v
       |Q_AC|
          |
          v
       Phase-C G_P
          |
          v
         P_HVAC
          |
          +-------------> MPC energy / price / grid objective
\end{verbatim}

\section{Empirically frozen mass-flow realization}

\subsection{Causal screening result}

The directed actuator experiments established that:
\begin{itemize}[leftmargin=2em]
\item the terminal/direct-air \texttt{Mass Flow Rate} actuator was writable but
did not causally control the realized supply flow in this PSZ topology;
\item the generic supply-node \texttt{Temperature Setpoint} actuator produced
no meaningful causal effect because the parent unitary systems operate in
\texttt{Load} control;
\item the \texttt{Fan Air Mass Flow Rate} actuator directly controlled the
realized supply flow;
\item unitary sensible-load request and DX-speed commands affected thermal
behavior and were retained for thermal-realization identification.
\end{itemize}

\subsection{Frozen fan mapping}

For the present two-mode PSZ/Fan:SystemModel implementation, C2.2 established
over Dining and Kitchen, heating and cooling, and the physical range
$0.2$--$1.0$ of design mass flow:
\begin{equation}
\boxed{
\mdot_{i}^{\mathrm{EP}} = 2 a_{m,i}
}
\qquad i\in\{D,K\}.
\label{eq:exact-fan-forward}
\end{equation}

The corresponding inverse is
\begin{equation}
\boxed{
a_{m,i}
=
\frac{1}{2}\mdot_i^\star .
}
\label{eq:exact-fan-inverse}
\end{equation}

This relation remained exact when thermal actuators were simultaneously active
in C2.3 and C2.4.  In C2.4 the maximum measured absolute flow error and maximum
fraction-of-design flow error were both reported as zero:
\begin{equation}
\max |e_i^{\mdot}| = 0,
\qquad
\max
\frac{|e_i^{\mdot}|}{\mdot_{i,\max}}
=0.
\label{eq:c24-flow-zero-error}
\end{equation}

Therefore Eq.~\eqref{eq:exact-fan-inverse} is \Frozen\ for this exact
RestaurantFastFood EnergyPlus 24.1 PSZ implementation.  It is a
plant-specific realization law, not a universal EnergyPlus identity.

\section{Why supply temperature requires a state-dependent realization map}

Changing mass flow alone changes supply temperature because fan flow, outdoor
air, coil operation, and unitary control are coupled.  Therefore the true
thermal realization cannot generally be written as a completely independent
map $T_{\mathrm{sa}}^\star\mapsto a_q$.

For the air path
\begin{equation}
T_{\mathrm{mix}}
\longrightarrow
T_{\mathrm{coil,out}}
\longrightarrow
T_{\mathrm{heat,out}}
\longrightarrow
T_{\mathrm{sa}},
\end{equation}
the zone-level heat can be decomposed as
\begin{align}
Q_{X,i}
&=
\mdot_i c_p
\left(
T_{\mathrm{sa},i}-T_{z,i}
\right)
\nonumber\\
&=
\mdot_i c_p
\left(
T_{\mathrm{mix},i}-T_{z,i}
\right)
+
\mdot_i c_p
\left(
T_{\mathrm{sa},i}-T_{\mathrm{mix},i}
\right).
\label{eq:airpath-decomposition}
\end{align}

Define
\begin{equation}
Q_{\mathrm{burden},i}
=
\mdot_i c_p
\left(
T_{\mathrm{mix},i}-T_{z,i}
\right),
\label{eq:airpath-burden}
\end{equation}
and
\begin{equation}
Q_{\mathrm{equip},i}
=
\mdot_i c_p
\left(
T_{\mathrm{sa},i}-T_{\mathrm{mix},i}
\right).
\label{eq:airpath-equipment}
\end{equation}
Then
\begin{equation}
Q_{X,i}
=
Q_{\mathrm{burden},i}
+
Q_{\mathrm{equip},i}.
\label{eq:airpath-sum}
\end{equation}

This decomposition explains why a heating coil can be strongly active while
the zone-level quantity $Q_{X,i}$ remains negative: the equipment may first be
overcoming a large negative outdoor/makeup-air burden.

\section{First-principles thermal-actuator candidate}

The simplest possible unitary sensible-load request associated with a desired
physical pair is
\begin{equation}
a_{q,i}^{(0)}
=
Q_{X,i}^\star
=
\mdot_i^\star c_p
\left(
T_{\mathrm{sa},i}^\star-T_{z,i}
\right).
\label{eq:naive-load-request}
\end{equation}

C2.3 showed that this direct physics request is already accurate over important
operating points for Dining heating, Dining cooling, and Kitchen cooling, but
is inadequate for Kitchen heating under a large makeup-air burden.  Therefore
Eq.~\eqref{eq:naive-load-request} is a useful feedforward term but is not yet a
universal low-level inverse.

\section{Local affine realization model}

A physically interpretable local model is
\begin{equation}
\boxed{
Q_{X,i}^{\mathrm{EP}}
=
\alpha_i(\xi_i)\,a_{q,i}
+
\beta_i(\xi_i)
+
\varepsilon_{Q,i},
}
\label{eq:local-affine}
\end{equation}
where:
\begin{itemize}[leftmargin=2em]
\item $\alpha_i(\xi_i)$ is the local effectiveness/gain of the unitary
sensible-load request;
\item $\beta_i(\xi_i)$ is the local zone-interface heat that exists at zero
low-level load request under the current air-loop/controller state;
\item $\varepsilon_{Q,i}$ collects local nonlinearity, mode switching, and
unmodeled effects.
\end{itemize}

When $\alpha_i\neq 0$ and the same local controller mode remains active, the
corresponding inverse candidate is
\begin{equation}
\boxed{
a_{q,i}^\star
=
\frac{
Q_{X,i}^\star-\beta_i(\xi_i)
}{
\alpha_i(\xi_i)
}
=
\frac{
\mdot_i^\star c_p(T_{\mathrm{sa},i}^\star-T_{z,i})
-\beta_i(\xi_i)
}{
\alpha_i(\xi_i)
}.
}
\label{eq:affine-inverse}
\end{equation}

Equation~\eqref{eq:affine-inverse} is \Candidate, not yet globally frozen.

\section{C2.4 counterfactual identification design}

C2.4 removed the sequential-pulse confound by running independent
counterfactual EnergyPlus simulations from the same historical trajectory and
the same selected operating instant.

\begin{verbatim}
counterfactual A: same history -> same state -> fixed m_dot -> Q_req = 0
counterfactual B: same history -> same state -> fixed m_dot -> Q_req = level 1
counterfactual C: same history -> same state -> fixed m_dot -> Q_req = level 2
counterfactual D: same history -> same state -> fixed m_dot -> Q_req = level 3
\end{verbatim}

For heating:
\begin{equation}
\Delta T^\star \in \{0,+4,+8,+12\}\ ^\circ\mathrm C,
\end{equation}
and for cooling:
\begin{equation}
\Delta T^\star \in \{0,-4,-8,-12\}\ ^\circ\mathrm C.
\end{equation}
The applied load request was
\begin{equation}
a_q
=
\mdot^\star c_p \Delta T^\star.
\label{eq:c24-command}
\end{equation}

The fixed physical flow was
\begin{equation}
\mdot_i^\star
=
0.66\,\mdot_{i,\max},
\label{eq:c24-flowtarget}
\end{equation}
implemented using Eq.~\eqref{eq:exact-fan-inverse}.

\section{C2.4 numerical results}

The C2.4 audit passed under the accepted EnergyPlus 24.1 candidate.  The
portable evidence bundle used for this document was
\begin{verbatim}
PINODE_EPSR_Day7C2_4_CounterfactualLoadMap_PASS_20260829_103409.zip

SHA256:
cbe597daa377d35ca85ca764dc34bb40d73be9fe014c993af7ccde9934ffb6f7
\end{verbatim}

\subsection{Primary affine fits reported by C2.4}

\begin{table}[htbp]
\centering
\small
\begin{tabular}{llrrrrl}
\toprule
Zone & Mode & $\alpha$ & $\beta$ [W] & $R^2$ & RMSE [W] & Classification \\
\midrule
Dining & Heating & 0.9610 & 292.1 & 0.99948 & 118.9 & strong affine \\
Dining & Cooling & 0.9128 & -776.7 & 0.99968 & 87.8 & strong affine \\
Kitchen & Heating & 0.1867 & -10514.8 & 0.28818 & 1295.1 & mode dependent \\
Kitchen & Cooling & 0.6331 & -3338.4 & 0.86299 & 1113.6 & mode dependent \\
\bottomrule
\end{tabular}
\caption{Primary C2.4 local affine fits using all four counterfactual points,
including the explicitly forced zero-load point.}
\label{tab:c24-affine}
\end{table}

Thus, Dining heating and cooling are strong local affine inverse candidates.
Kitchen cannot yet be represented by one global affine line across the
zero-load and active-actuation points.

\subsection{Counterfactual points}

\begin{table}[htbp]
\centering
\small
\begin{tabular}{llrrr}
\toprule
Zone & Mode & $\Delta T^\star$ [$^\circ$C] &
$a_q$ [W] & $Q_X^{\mathrm{EP}}$ [W] \\
\midrule
Dining & Heating & 0 & 0 & 241.1 \\
Dining & Heating & +4 & 4842.6 & 5103.3 \\
Dining & Heating & +8 & 9685.2 & 9438.2 \\
Dining & Heating & +12 & 14527.8 & 14309.3 \\
\midrule
Dining & Cooling & 0 & 0 & -863.4 \\
Dining & Cooling & -4 & -4842.6 & -5112.4 \\
Dining & Cooling & -8 & -9685.2 & -9526.3 \\
Dining & Cooling & -12 & -14527.8 & -14126.4 \\
\midrule
Kitchen & Heating & 0 & 0 & -9097.9 \\
Kitchen & Heating & +4 & 3948.2 & -11720.0 \\
Kitchen & Heating & +8 & 7896.5 & -9407.1 \\
Kitchen & Heating & +12 & 11844.7 & -7412.1 \\
\midrule
Kitchen & Cooling & 0 & 0 & -4557.8 \\
Kitchen & Cooling & -4 & -3948.2 & -4187.7 \\
Kitchen & Cooling & -8 & -7896.5 & -7980.5 \\
Kitchen & Cooling & -12 & -11844.7 & -11626.0 \\
\bottomrule
\end{tabular}
\caption{C2.4 settled zone-interface heat for the independently restarted
counterfactual simulations.}
\label{tab:c24-points}
\end{table}

\subsection{Supply-temperature tracking}

The settled supply-temperature tracking errors
$T_{\mathrm{sa}}^{\mathrm{EP}}-T_{\mathrm{sa}}^\star$ were approximately:
\begin{align}
\text{Dining heating:}\quad
&\{+0.199,+0.215,-0.204,-0.180\}\ ^\circ\mathrm C,
\\
\text{Dining cooling:}\quad
&\{-0.713,-0.223,+0.131,+0.332\}\ ^\circ\mathrm C,
\\
\text{Kitchen heating:}\quad
&\{-9.217,-15.874,-17.530,-19.509\}\ ^\circ\mathrm C,
\\
\text{Kitchen cooling:}\quad
&\{-4.618,-0.243,-0.085,+0.222\}\ ^\circ\mathrm C.
\end{align}

The Kitchen zero-load cooling point and all Kitchen heating points clearly
belong to controller/air-path regimes that cannot be captured by one smooth
global affine map centered at zero load.

\subsection{Active-regime post-hoc fits}

For interpretation, fitting only the three nonzero active-command points gives:
\begin{align}
\text{Dining heating:}\quad
&Q_X^{\mathrm{EP}}
\approx
0.9505\,a_q + 410.9,
\qquad R^2\approx0.99887,
\\
\text{Dining cooling:}\quad
&Q_X^{\mathrm{EP}}
\approx
0.9307\,a_q - 574.4,
\qquad R^2\approx0.99986,
\\
\text{Kitchen cooling:}\quad
&Q_X^{\mathrm{EP}}
\approx
0.9420\,a_q - 493.1,
\qquad R^2\approx0.99987,
\\
\text{Kitchen heating:}\quad
&Q_X^{\mathrm{EP}}
\approx
0.5455\,a_q - 13821.0,
\qquad R^2\approx0.99819.
\end{align}

These active-only regressions are derived diagnostics from the C2.4
counterfactual points, not yet frozen paper parameters.  They suggest a
piecewise/controller-mode realization structure:
\begin{equation}
Q_X^{\mathrm{EP}}
=
\alpha_{i,\sigma} a_q
+
\beta_{i,\sigma}
+
\varepsilon,
\label{eq:piecewise-affine}
\end{equation}
where $\sigma$ denotes an identified HVAC/controller mode.

The Kitchen-heating active-region result is particularly important: it can be
locally affine while still having a very large negative offset.  Thus a
positive low-level heating request does not imply that the resulting zone-level
heat is positive unless the request first exceeds the current air-path burden.

\section{Candidate realization-state vector}

The realization context should be kept as small and physically interpretable as
the data permit.  A candidate vector is
\begin{equation}
\xi_i
=
\begin{bmatrix}
T_{z,i} &
T_{\mathrm{mix},i} &
W_{\mathrm{mix},i} &
T_{\mathrm{oa}} &
\mdot_i &
\sigma_i &
r_i
\end{bmatrix}^{\top},
\label{eq:xi-candidate}
\end{equation}
where $r_i$ may collect only additional equipment/controller states shown by
future identification to be necessary.  \OpenItem\ The final components of
$\xi_i$ must be selected from directed numerical evidence rather than included
by default.

A preferred hierarchy is:
\begin{enumerate}[leftmargin=2em]
\item exact analytical relation where available;
\item physics-based feedforward plus calibrated gain/offset;
\item piecewise-affine or monotone interpolation by controller mode;
\item compact nonlinear regression only if simpler structures fail held-out
validation;
\item constrained low-level optimization when the inverse is nonunique or the
requested pair is infeasible.
\end{enumerate}

\section{General nonlinear realizer}

The most general forward map is
\begin{equation}
u_i^{\mathrm{EP}}
=
g_i(a_i,\xi_i).
\label{eq:general-forward}
\end{equation}

If a closed-form inverse is unavailable, define the realizer by
\begin{equation}
\boxed{
a_{i,k}^{\star}
=
\arg\min_{a\in\mathcal A_i}
\left\|
g_i(a,\xi_{i,k})
-
u_{i,k}^{\star}
\right\|_{W_u}^{2}
+
\lambda_a
\left\|
a-a_{i,k-1}
\right\|_{W_a}^{2}.
}
\label{eq:low-level-realizer-optimization}
\end{equation}

Optional additional penalties can include low-level electrical power or
actuator wear, but they should not duplicate objectives already present in the
upper MPC without a documented reason.

For the current plant, the flow component can be removed from the numerical
inverse because Eq.~\eqref{eq:exact-fan-inverse} is exact:
\begin{equation}
a_{m,i}^{\star}
=
\frac12 \mdot_i^\star.
\end{equation}
The remaining low-level problem can therefore focus primarily on thermal
realization.

\section{Realizable physical-input set}

For a fixed HVAC state/context $\xi_i$, define
\begin{equation}
\mathcal U_i^{\mathrm{EP}}(\xi_i)
=
\left\{
u:
\exists a\in\mathcal A_i
\ \mathrm{s.t.}\
u=g_i(a,\xi_i)
\right\}.
\label{eq:ep-feasible-set}
\end{equation}

The thermal model should not be driven by control requests far outside the
region seen during identification/training.  Therefore define
\begin{equation}
\boxed{
\mathcal U_i^{\mathrm{MPC}}(\xi_i)
=
\mathcal U_i^{\mathrm{train}}
\cap
\mathcal U_i^{\mathrm{EP}}(\xi_i).
}
\label{eq:mpc-feasible-set}
\end{equation}

The upper MPC must impose
\begin{equation}
u_{i,k+j|k}^\star
\in
\mathcal U_i^{\mathrm{MPC}}
\left(
\xi_{i,k+j|k}
\right).
\label{eq:mpc-realizable-constraint}
\end{equation}

If a requested pair is outside the realizable set, an optional projection or
low-level constrained realizer may be used:
\begin{equation}
\bar u_i^\star
=
\operatorname*{arg\,min}_{u\in\mathcal U_i^{\mathrm{EP}}(\xi_i)}
\|u-u_i^\star\|_{W_u}^{2}.
\label{eq:realizable-projection}
\end{equation}
Any use of Eq.~\eqref{eq:realizable-projection} must be logged because it means
the plant did not execute the MPC request exactly.

\section{Actuator-realization error}

Define
\begin{equation}
e_{i,k}^{u}
=
u_{i,k}^{\mathrm{EP}}
-
u_{i,k}^{\star}
=
\begin{bmatrix}
e_{i,k}^{\mdot}\\
e_{i,k}^{T}
\end{bmatrix},
\label{eq:realization-error-vector}
\end{equation}
where
\begin{align}
e_{i,k}^{\mdot}
&=
\mdot_{i,k}^{\mathrm{EP}}
-
\mdot_{i,k}^{\star},
\\
e_{i,k}^{T}
&=
T_{\mathrm{sa},i,k}^{\mathrm{EP}}
-
T_{\mathrm{sa},i,k}^{\star}.
\end{align}

The exact raw air-side heat error is
\begin{align}
e_{Q_X,i}
&=
Q_{X,i}^{\mathrm{EP}}
-
Q_{X,i}^{\star}
\nonumber\\
&=
c_p
\left[
e_i^{\mdot}
\left(
T_{\mathrm{sa},i}^\star-T_{z,i}
\right)
+
\mdot_i^{\mathrm{EP}}e_i^T
\right].
\label{eq:exact-qx-error}
\end{align}

Hence
\begin{equation}
\boxed{
|e_{Q_X,i}|
\le
c_p
\left(
|e_i^{\mdot}|
|T_{\mathrm{sa},i}^\star-T_{z,i}|
+
|\mdot_i^{\mathrm{EP}}|
|e_i^T|
\right).
}
\label{eq:qx-error-bound}
\end{equation}

Because the current flow realization is exact in the tested plant,
Eq.~\eqref{eq:qx-error-bound} reduces experimentally to
\begin{equation}
|e_{Q_X,i}|
=
c_p
|\mdot_i^\star|
|e_i^T|
\end{equation}
for those validated regimes, up to numerical tolerance.

If $G_Q$ is locally Lipschitz with constant $L_Q$ over the operating region,
then
\begin{equation}
|e_{Q_{\mathrm{AC}},i}|
\le
L_Q |e_{Q_X,i}|.
\label{eq:qac-error-lipschitz}
\end{equation}
If $G_P$ is locally Lipschitz with constant $L_P$, then a corresponding
conditional electrical-power bound is
\begin{equation}
|e_{P_{\mathrm{HVAC}},i}|
\le
L_P |e_{Q_{\mathrm{AC}},i}|.
\label{eq:phvac-error-lipschitz}
\end{equation}
Equations~\eqref{eq:qac-error-lipschitz} and
\eqref{eq:phvac-error-lipschitz} are diagnostic bounds conditional on the
stated local regularity and are not MPC stability theorems.

\section{Continuous-time thermal prediction inside MPC}

For each frozen thermal-model family, retain the continuous-time dynamics
\begin{equation}
\dot x(t)
=
F_\theta
\left(
x(t),
\Qac(t),
d(t)
\right).
\label{eq:continuous-model}
\end{equation}

For MPC step $k$ and candidate physical control
$u_{k+j|k}^{\star}$:
\begin{enumerate}[leftmargin=2em]
\item compute the desired raw air-side heat using
Eq.~\eqref{eq:desired-airside-heat};
\item apply the frozen Phase-C $G_Q$ mapping;
\item numerically integrate the continuous-time model over
$\Delta t_{\mathrm{MPC}}=300$ s;
\item compute $\Phvac$ separately using $G_P$ for the objective.
\end{enumerate}

Using a fixed-step RK4 operator $\Phi_{\mathrm{RK4}}$,
\begin{equation}
x_{k+1}
=
\Phi_{\mathrm{RK4}}
\left(
x_k,
Q_{\mathrm{AC},k},
d_k;
\Delta t_{\mathrm{MPC}}
\right).
\label{eq:rk4-step}
\end{equation}

\section{Observer interface}

This realization document does not replace the Part-6 observer definitions.
The MPC should receive the state estimate
\begin{equation}
\hat x_k
=
\mathcal O
\left(
y_{0:k},
u_{0:k-1},
d_{0:k}
\right).
\label{eq:generic-observer}
\end{equation}

For learned two-capacitance models, the measured zone temperature directly
anchors the observed state and the frozen causal encoder estimates only hidden
states.  For Inverse PINN-RC two-capacitance cases, the Part-6
Luenberger/Kalman-family observer options remain applicable.  The low-level
realizer may consume measured EnergyPlus air-path quantities in $\xi_k$ without
changing the thermal observer definition.

\section{Upper-level MPC formulation}

Let the prediction horizon be $N$.  A representative scalarized objective is
\begin{equation}
\min_{\{u^\star_{k+j|k}\}_{j=0}^{N-1}}
J_k
=
w_T J_T
+
w_E J_E
+
w_C J_C
+
w_G J_G
+
w_{\Delta u}J_{\Delta u}
+
w_s J_s
+
w_R J_R,
\label{eq:mpc-objective}
\end{equation}
where:
\begin{align}
J_T
&=
\sum_{j=1}^{N}
\ell_T
\left(
T_{z,k+j|k},
T_{z,k+j}^{\min},
T_{z,k+j}^{\max}
\right),
\\
J_E
&=
\sum_{j=0}^{N-1}
P_{\mathrm{HVAC},k+j|k}\Delta t_{\mathrm{MPC}},
\\
J_C
&=
\sum_{j=0}^{N-1}
c_{k+j}
P_{\mathrm{HVAC},k+j|k}
\Delta t_{\mathrm{MPC}},
\\
J_G
&=
\sum_{j=0}^{N-1}
\left(
P_{\mathrm{HVAC},k+j|k}
-
P_{k+j}^{\mathrm{ref}}
\right)^2,
\\
J_{\Delta u}
&=
\sum_{j=0}^{N-1}
\left\|
u^\star_{k+j|k}
-
u^\star_{k+j-1|k}
\right\|_{R}^{2},
\\
J_R
&=
\sum_{j=0}^{N-1}
\widehat{\ell}_{R}
\left(
u^\star_{k+j|k},\xi_{k+j|k}
\right).
\end{align}

$J_R$ is optional and may penalize requests predicted to have large
actuator-realization error or poor local conditioning.  It should be included
only after its empirical benefit is demonstrated.

The principal constraints are
\begin{align}
x_{k+j+1|k}
&=
\Phi_{\mathrm{RK4}}
\left(
x_{k+j|k},
Q_{\mathrm{AC},k+j|k},
d_{k+j|k};
\Delta t_{\mathrm{MPC}}
\right),
\\
Q_{X,k+j|k}^{\star}
&=
\mdot_{k+j|k}^{\star}
c_p
\left(
T_{\mathrm{sa},k+j|k}^{\star}
-
T_{z,k+j|k}
\right),
\\
Q_{\mathrm{AC},k+j|k}
&=
G_Q
\left(
Q_{X,k+j|k}^{\star},
\zeta_{Q,k+j|k}
\right),
\\
P_{\mathrm{HVAC},k+j|k}
&=
G_P
\left(
|Q_{\mathrm{AC},k+j|k}|,
\zeta_{P,k+j|k}
\right),
\\
u_{k+j|k}^\star
&\in
\mathcal U^{\mathrm{MPC}}
\left(
\xi_{k+j|k}
\right).
\end{align}

No binary heat/cool optimization variables are introduced unless future
experiments prove them necessary.  The existing continuous auxiliary/epigraph
treatment of $|\Qac|$ should be retained for the NLP.

\section{Identity two-zone control}

For the identity RestaurantFastFood representation,
\begin{equation}
u_k^\star
=
\begin{bmatrix}
u_{D,k}^{\star}\\
u_{K,k}^{\star}
\end{bmatrix}
=
\begin{bmatrix}
\mdot_{D,k}^\star\\
T_{\mathrm{sa},D,k}^\star\\
\mdot_{K,k}^\star\\
T_{\mathrm{sa},K,k}^\star
\end{bmatrix}.
\label{eq:identity-control-vector}
\end{equation}

The two low-level realizers are allowed to have different thermal maps:
\begin{align}
a_{D,k}
&=
f_D(u_{D,k}^\star,\xi_{D,k}),
\\
a_{K,k}
&=
f_K(u_{K,k}^\star,\xi_{K,k}).
\end{align}
C2.4 strongly supports keeping these maps zone-specific rather than forcing one
shared thermal inverse.

\section{All-to-one aggregate-model control and allocation}

When the predictive thermal model is all-to-one, the predictive state is
aggregate but the EnergyPlus plant still contains two physical PSZ systems.
Therefore an explicit physical-input allocation layer is required.

Let the aggregate MPC request be
\begin{equation}
u_{a,k}^{\star}
=
\begin{bmatrix}
\mdot_{a,k}^{\star}\\
T_{\mathrm{sa},a,k}^{\star}
\end{bmatrix}.
\end{equation}

Define a detailed allocation map
\begin{equation}
u_{k}^{z,\star}
=
\mathcal A_u
\left(
u_{a,k}^{\star},
\xi_{D,k},
\xi_{K,k}
\right)
=
\begin{bmatrix}
u_{D,k}^{\star}\\
u_{K,k}^{\star}
\end{bmatrix}.
\label{eq:aggregate-allocation}
\end{equation}

The implemented allocation must be checked against the aggregation input map
$P_u$.  Define the allocation residual
\begin{equation}
r_{u,k}
=
u_{a,k}^{\star}
-
P_u u_{k}^{z,\star}.
\label{eq:allocation-residual}
\end{equation}
This quantity connects directly to the existing aggregation-theorem diagnostics:
an aggregate thermal request need not correspond to one unique detailed
Dining/Kitchen actuator allocation.

\OpenItem\ The final $\mathcal A_u$ must be selected from the same actuator
realizability evidence and should not be assumed to be equal splitting merely
because the thermal aggregation uses equal temperature weights.

\section{Closed-loop execution algorithm}

At each real EnergyPlus control instant $k$:

\begin{enumerate}[leftmargin=2em]
\item Read measurements and update the thermal state estimate $\hat x_k$.
\item Build weather, disturbance, price, and grid/reference forecasts.
\item Solve the upper-level MPC in physical coordinates
$(\mdot^\star,T_{\mathrm{sa}}^\star)$.
\item For an aggregate predictive model, allocate the aggregate physical request
to Dining/Kitchen requests using $\mathcal A_u$.
\item For each physical zone, evaluate the low-level realization context
$\xi_{i,k}$.
\item Apply the exact fan inverse
$a_{m,i,k}=\tfrac12\mdot_{i,k}^{\star}$.
\item Evaluate the thermal realizer $a_{q,i,k}$ and, only if justified,
$a_{s,i,k}$.
\item Apply low-level commands through the EnergyPlus API.
\item Measure
$\mdot_{i,k}^{\mathrm{EP}}$ and
$T_{\mathrm{sa},i,k}^{\mathrm{EP}}$.
\item Compute and log $e^u_{i,k}$,
$e_{Q_X,i,k}$, realized $\Qac$, and realized $\Phvac$.
\item Advance EnergyPlus to the next 300-s MPC instant and repeat.
\end{enumerate}

\begin{verbatim}
MEASURE
  |
  v
OBSERVER / ENCODER
  |
  v
MPC in physical coordinates
  |
  |  (m_dot*, T_sa*)
  v
REALIZABILITY / ALLOCATION
  |
  v
LOW-LEVEL REALIZER
  |
  +--> a_fan = 0.5 m_dot*
  |
  +--> a_q = thermal inverse(...)
  |
  +--> optional DX command
  |
  v
ENERGYPLUS
  |
  v
(m_dot_EP, T_sa_EP)
  |
  +--> realization-error metrics
  |
  +--> Q_HVAC,X -> G_Q -> Q_AC -> thermal state
  |
  +--> |Q_AC| -> G_P -> P_HVAC
  |
  `----------------------------------> next MPC step
\end{verbatim}

\section{Required identification data}

Future actuator-realization experiments should retain, at minimum:

\subsection*{Zone/MPC-level signals}
\begin{equation}
\left\{
T_z,\;
\mdot^{\star},\;
T_{\mathrm{sa}}^{\star},\;
\mdot^{\mathrm{EP}},\;
T_{\mathrm{sa}}^{\mathrm{EP}},\;
Q_X^{\star},\;
Q_X^{\mathrm{EP}}
\right\}.
\end{equation}

\subsection*{Low-level actuator signals}
\begin{equation}
\left\{
a_m,\;
a_q,\;
a_s,\;
\text{fan power},\;
Q_{\mathrm{heat,coil}},\;
Q_{\mathrm{cool,coil}}
\right\}.
\end{equation}

\subsection*{Air-path signals}
At the return, mixed-air, cooling-coil outlet, heating-coil outlet, PSZ supply
outlet, and zone direct-air inlet:
\begin{equation}
\left\{
T,\;
\mdot,\;
W
\right\}.
\end{equation}

\subsection*{Context/mode signals}
Retain only quantities needed to explain mode dependence, such as:
\begin{equation}
\left\{
T_{\mathrm{oa}},
T_{\mathrm{mix}},
W_{\mathrm{mix}},
\text{heating/cooling/native controller mode},
\text{unitary/DX state}
\right\}.
\end{equation}

All identification summaries must be time weighted when EnergyPlus changes the
system-timestep duration during HVAC iteration.

\section{Identification roadmap after C2.4}

The current evidence suggests the following directed sequence.

\subsection*{C2.5: controller-mode segmentation}
Determine whether Kitchen discontinuities are explained by a small number of
discrete native controller regimes.  Candidate segmentation variables should
be inferred from the actual air-loop/controller diagnostics, not guessed.

\subsection*{C2.6: state-conditioned DOE}
For each physically relevant mode, vary:
\begin{equation}
\mdot/\mdot_{\max},
\qquad
T_z,
\qquad
T_{\mathrm{mix}}\ \text{or}\ T_{\mathrm{oa}},
\qquad
a_q,
\end{equation}
using independent counterfactuals.  Estimate
\begin{equation}
\alpha_i(\xi),\qquad \beta_i(\xi),
\end{equation}
and determine whether a piecewise-affine model is sufficient.

\subsection*{C2.7: held-out inverse validation}
Given desired $(\mdot^\star,T_{\mathrm{sa}}^\star)$ pairs not used to fit the realizer,
compute the low-level command and evaluate:
\begin{align}
\mathrm{MAE}_{\mdot},
\quad
\mathrm{RMSE}_{\mdot},
\quad
\max |e^{\mdot}|,
\\
\mathrm{MAE}_{T_{\mathrm{sa}}},
\quad
\mathrm{RMSE}_{T_{\mathrm{sa}}},
\quad
\max |e^T|.
\end{align}

\subsection*{C2.8: feasible-region identification}
Map the state-dependent realizable set
$\mathcal U_i^{\mathrm{EP}}(\xi)$ and intersect it with the TRAIN-supported
control region.

\subsection*{C2.9: frozen actuator contract}
Freeze exact EnergyPlus actuator keys, units, call points, scaling, clipping,
controller-mode logic, inverse model, and fallback behavior.

\section{Status matrix}

\begin{longtable}{p{0.27\textwidth}p{0.20\textwidth}p{0.44\textwidth}}
\toprule
Item & Status & Current conclusion \\
\midrule
\endfirsthead
\toprule
Item & Status & Current conclusion \\
\midrule
\endhead
EnergyPlus 24.1 baseline & Frozen & Unchanged SHA-verified transition candidate accepted after airflow, DHW, and DX warning forensics. \\
MPC physical variables & Locked & Upper MPC remains in $(\mdot,T_{\mathrm{sa}})$ physical coordinates. \\
Terminal mass-flow actuator & Rejected & Writable but did not causally control actual supply flow in this topology. \\
Supply-node temperature setpoint & Rejected & No causal effect under parent unitary \texttt{Load} control. \\
Fan mass-flow actuator & Frozen & Exact relation $\mdot^{EP}=2a_m$ over tested range/modes. \\
Fan inverse & Frozen & $a_m=\tfrac12\mdot^\star$. \\
Unitary sensible-load actuator & Retained & Strong thermal-control candidate; direct physical request works well in several active regimes. \\
DX-speed actuator & Retained for study & Potential cooling-side control coordinate; not yet selected as final thermal realizer. \\
Dining thermal inverse & Strong candidate & Local affine behavior in C2.4 with $R^2>0.999$ heating and cooling. \\
Kitchen cooling inverse & Piecewise candidate & Active nonzero cooling points are nearly affine, but zero-load/native-controller transition breaks one global line. \\
Kitchen heating inverse & Open & Large makeup-air burden and controller-mode dependence require state/mode-conditioned realization. \\
Realizable $(\mdot,T_{\mathrm{sa}})$ set & Open & Must be identified numerically and intersected with training support. \\
Aggregate-to-zone allocator & Open & Must preserve aggregate meaning while respecting different Dining/Kitchen realizability. \\
Final closed-loop MPC & Pending & Implement only after actuator realizer and feasible set are frozen. \\
\bottomrule
\end{longtable}

\section{Paper-facing interpretation}

The intended paper statement is not that EnergyPlus exposes ideal independent
$\mdot$ and $T_{\mathrm{sa}}$ actuators.  Rather, the control architecture is:
\begin{quote}
The MPC optimizes physically interpretable supply-air mass-flow and
supply-temperature trajectories.  A separately identified actuator-realization
layer maps these physical requests to the low-level commands of the native
EnergyPlus packaged-unitary systems.  The map is calibrated by directed
counterfactual simulation, constrained to the empirically realizable operating
region, and validated by commanded-versus-realized airflow and supply
temperature.
\end{quote}

This formulation preserves the physical meaning of the thermal model and the
MPC while treating equipment implementation as a plant-specific realization
problem.

\section{Claims that should not be made}

The current evidence does \emph{not} support the following universal claims:
\begin{itemize}[leftmargin=2em]
\item ``EnergyPlus directly controls arbitrary $\mdot$ and $T_{\mathrm{sa}}$ pairs.''
\item ``The fan factor of two is universal across EnergyPlus fan systems.''
\item ``One affine thermal inverse is valid for all zones and all controller
modes.''
\item ``Positive heating-coil power implies positive heat delivery relative to
the zone.''
\item ``The zero-load point belongs to the same local controller regime as
active heating/cooling.''
\item ``Actuator realizability proves MPC recursive feasibility or closed-loop
stability.''
\item ``$\Qac$ and $\Phvac$ are interchangeable quantities.''
\end{itemize}

\section{Provenance}

The persistent scientific root for this work is
\begin{verbatim}
C:\Users\ninad\Dropbox\NinadGaikwad_PhD\Gaikwad_Research\
From_WSU_OneDrive\BuildingModelingProject_Condensed\
Data\ScaleBridge\Paper_PINODE_EPSR
\end{verbatim}

Relevant actuator-audit families are stored below
\begin{verbatim}
03_mpc_actuator_audit\
\end{verbatim}
with unique timestamped folders.

This document incorporates numerical results from C1, C2.1, C2.2, C2.3, and
the C2.4 PASS bundle identified above.  It is designed to live alongside the
existing mathematical contracts under
\begin{verbatim}
Paper_PINODE_EPSR\docs\mathematics\contracts\
\end{verbatim}
as a new actuator-realization extension to the existing Part-6 MPC/observer
contract.

\section{Recommended repository filename}

A suitable repository filename is
\begin{verbatim}
PINODE_EPSR_Part7_MPC_Actuator_Realization_Detailed.tex
\end{verbatim}
or, if Part 6 is intentionally kept as the sole MPC document, this material
can be merged into a versioned extension of Part 6 after the remaining
Kitchen-mode identification is frozen.

\section{Current mathematical endpoint}

The present experimentally supported architecture can be summarized as
\begin{equation}
\boxed{
u_{i,k}^{\star}
=
\begin{bmatrix}
\mdot_{i,k}^{\star}\\
T_{\mathrm{sa},i,k}^{\star}
\end{bmatrix}
}
\end{equation}
from the upper MPC, followed by
\begin{equation}
\boxed{
a_{m,i,k}
=
\frac12\mdot_{i,k}^{\star}
}
\end{equation}
and a thermal realizer
\begin{equation}
\boxed{
a_{q,i,k}
=
f_{q,i}
\left(
\mdot_{i,k}^{\star},
T_{\mathrm{sa},i,k}^{\star},
\xi_{i,k}
\right),
}
\end{equation}
with the physics-informed piecewise-affine candidate
\begin{equation}
\boxed{
a_{q,i,k}
\approx
\frac{
\mdot_{i,k}^{\star}c_p
(T_{\mathrm{sa},i,k}^{\star}-T_{z,i,k})
-
\beta_{i,\sigma}(\xi_{i,k})
}{
\alpha_{i,\sigma}(\xi_{i,k})
}.
}
\label{eq:current-endpoint}
\end{equation}

The resulting realized physical action is
\begin{equation}
u_{i,k}^{\mathrm{EP}}
=
g_i
\left(
a_{i,k},\xi_{i,k}
\right),
\end{equation}
and the closed-loop experiment explicitly evaluates
\begin{equation}
u_{i,k}^{\mathrm{EP}}
-
u_{i,k}^{\star}.
\end{equation}

Equation~\eqref{eq:current-endpoint} is the present mathematical bridge between
the physically interpretable MPC variables and the native EnergyPlus PSZ
actuator system.  The remaining directed numerical work is to identify the
smallest state/mode dependence required for $\alpha$ and $\beta$, validate the
inverse on held-out counterfactuals, and freeze the realizable control set.

\end{document}
